\chapter{Physics-Informed Attention Biases}

Standard self-attention assigns relevance purely from learned query--key
inner products, treating all token pairs equally. In a multi-top event,
however, particles that share a decay vertex, are close in angular
separation, or interact via a specific SM coupling are physically more
related than arbitrary pairs. This chapter asks whether injecting such
structure as additive biases on the attention logits improves
classification, and what the learned bias weights reveal about how the
model organises particle interactions.

Four bias families are available: a Lorentz-scalar network mapping
pairwise invariants to bias scores~(B1-L), a learnable type-pair
interaction table optionally seeded from SM coupling structure~(B1-T),
a fixed SM interaction prior~(B1-S), and a MET-conditioned global
modulation~(B1-G). All families share a zero-initialised scalar gate,
so adding any bias leaves the baseline unchanged at initialisation.
All experiments in this chapter fix D02~=~\texttt{true}: MET is always
present, which is required for B1-G and motivated by the Ch.\,4 finding.

The primary lens is \textit{interpretability}: performance differences
are expected to be small, and the main contribution is understanding
what the models learn.

% ──────────────────────────────────────────────
\section{Bias family comparison (Exp.\ 5A)}

Compares each bias family individually and in full combination against
an unbiased baseline. Identifies which family, if any, provides a
consistent performance benefit, and whether the full stack beats the
sum of its parts.

\begin{table}[h]
\centering\small
\begin{tabular}{llr}
\toprule
Axis & Values & Config \\
\midrule
B1 — Bias selector & \texttt{none} & 6 \\
 & \texttt{lorentz\_scalar} & \\
 & \texttt{typepair\_kinematic} & \\
 & \texttt{sm\_interaction} & \\
 & \texttt{global\_conditioned} & \\
 & \texttt{lorentz+typepair+sm+global} & \\
Seeds & & 3 \\
\midrule
Total & & \textbf{18} \\
\bottomrule
\end{tabular}
\caption{Exp.\ 5A. All other axes at chapter baseline:
T1=\texttt{identity}, D02=\texttt{true}, E1=\texttt{sinusoidal}.}
\end{table}

Output: seed-averaged AUROC bar chart. Attention maps extracted at
inference time for the \texttt{none} baseline and full combination,
averaged over a fixed held-out batch and labelled by particle type.

% ──────────────────────────────────────────────
\section{Lorentz feature set and bias network (Exp.\ 5B)}

Isolates the Lorentz-scalar family to ask two questions: which pairwise
invariants are most informative (B1-L1), and whether replacing the
pointwise MLP with a KAN improves the learned response (B1-L2). Sparse
gating (B1-L5) adds per-feature sigmoid gates that directly expose
learned feature importance without ablation. The KAN activation curves
and converged gate values are extracted post-hoc from checkpoints ---
no re-training required.

\begin{table}[h]
\centering\small
\begin{tabular}{llr}
\toprule
Axis & Values & Config \\
\midrule
B1-L1 — Feature set & \texttt{[m2]} & 5 \\
 & \texttt{[deltaR]} & \\
 & \texttt{[m2, deltaR]} & \\
 & \texttt{[log\_kt, z, deltaR, log\_m2]} & \\
 & \texttt{[m2, deltaR, log\_m2, log\_kt, z, deltaR\_ptw]} & \\
B1-L2 — MLP type & \texttt{standard}, \texttt{kan} & 2 \\
B1-L5 — Sparse gating & \texttt{false}, \texttt{true} & 2 \\
Seeds & & 3 \\
\midrule
Total & & \textbf{60} \\
\bottomrule
\end{tabular}
\caption{Exp.\ 5B. B1=\texttt{lorentz\_scalar} throughout.}
\end{table}

% ──────────────────────────────────────────────
\section{Type-pair initialisation and freeze (Exp.\ 5C)}

The type-pair table is an $N_\text{type} \times N_\text{type}$ matrix
of learnable per-type-pair attention scores. This experiment asks
whether seeding it from SM coupling structure accelerates learning, and
whether a frozen physics prior alone is sufficient. The converged table
is visualised as a heatmap and compared across initialisation conditions.

\begin{table}[h]
\centering\small
\begin{tabular}{llr}
\toprule
Axis & Values & Config \\
\midrule
B1-T2 — Init & \texttt{none} & 5 \\
 & \texttt{binary} & \\
 & \texttt{binary} (frozen) & \\
 & \texttt{fixed\_coupling} & \\
 & \texttt{fixed\_coupling} (frozen) & \\
B1-T3 — Freeze & \texttt{false}, \texttt{true} & \\
Seeds & & 3 \\
\midrule
Total & & \textbf{15} \\
\bottomrule
\end{tabular}
\caption{Exp.\ 5C. Combination \texttt{init=none, freeze=true} excluded.
B1=\texttt{typepair\_kinematic} throughout.}
\end{table}

% ──────────────────────────────────────────────
\section{SM interaction mode (Exp.\ 5D)}

Tests whether increasing physical fidelity in the SM interaction
prior --- from a binary coupling mask through fixed coupling constants
to running couplings --- translates into measurable performance gains.
Expected to be a null result; feeds directly into Ch.\,8.\\

\begin{table}[h]
\centering\small
\begin{tabular}{llr}
\toprule
Axis & Values & Config \\
\midrule
B1-S1 — SM mode & \texttt{binary} & 3 \\
 & \texttt{fixed\_coupling} & \\
 & \texttt{running\_coupling} & \\
Seeds & & 3 \\
\midrule
Total & & \textbf{9} \\
\bottomrule
\end{tabular}
\caption{Exp.\ 5D. B1=\texttt{sm\_interaction} throughout.}
\end{table}


\subsection*{Chapter ran summary}

\begin{table}[h]
\centering\small
\begin{tabular}{llr}
\toprule
Experiment & Primary axes & Runs \\
\midrule
5A — Family comparison    & B1          & 18 \\
5B — Lorentz ablation     & B1-L1, B1-L2, B1-L5 & 60 \\
5C — Type-pair table      & B1-T2, B1-T3 & 15 \\
5D — SM interaction mode  & B1-S1        &  9 \\
\midrule
\textbf{Total} & & \textbf{102} \\
\bottomrule
\end{tabular}
\end{table}
